%% =====================================================================
%%  T-01-05  Results Are the Only Test
%%  -------------------------------------------------------------------
%%  One idea per file. Core is mandatory. Slots in canonical order.
%%  Max 700 words. No layout commands. No filler. New source material
%%  goes in sources/<BOOK>/T-01-05.tex -- never by rewriting this file.
%% =====================================================================
%% @kind: topic
%% @id: T-01-05
%% @title: Results Are the Only Test
%% @subtitle: The practitioner's epistemology, and where it fails
%% @chapter: c01
%% @order: 50
%% @status: stable
%% @sources: B1, B3
%% @updated: 2026-09-19

\topic{T-01-05}{Results Are the Only Test}{The practitioner's epistemology, and where it fails}

\begin{core}
Both sources judge a tactic by whether it worked, not by whether it is elegant, principled or well-argued. That standard is what makes the material usable and what makes it unreliable, and it needs to be held consciously rather than absorbed.
\end{core}
\begin{mechanism}
A results standard selects hard for techniques that exploit stable features of human behaviour --- loss aversion, need for consistency, status sensitivity, conflict avoidance --- because those pay out repeatedly across situations. It also selects against techniques that are merely correct. The weakness is structural: a single success is not evidence, and the practitioner who learns from outcomes is learning from a small, self-selected, survivorship-biased sample. The con artist's craft is genuinely well calibrated on the marks he meets often and badly calibrated on everything else. The discipline is worth adopting with the sample-size problem attached.
\end{mechanism}
\begin{tells}
  \item A method is defended by a story rather than by a mechanism.
  \item Failures are attributed to the situation; successes are attributed to the method.
  \item No one can say what would count as evidence that the method had stopped working.
\end{tells}
\begin{moves}
  \item Ask of any tactic: what is the mechanism, and would I expect it to work on a person who does not share this mark's disposition?
  \item Keep the failures. They are the only informative data in the set.
  \item Prefer tactics whose mechanism you can state, so that you can predict where they break.
\end{moves}
\begin{counters}
  \item When someone justifies a method by their success with it, ask what happened the times it did not work.
  \item Refuse the survivorship sample: the people who were taken in are not in the room to describe it.
  \item Judge a technique by its mechanism, not by its testimonial.
\end{counters}
\begin{cost}
Applied to yourself, a pure results standard will eventually license something indefensible that happened to work once. Applied to others' claims, it is excellent. The asymmetry is the point: demand mechanism from yourself and accept outcomes from nobody, including the sources of this book. Both are confident in ways their evidence does not support.
\end{cost}

\sources{B1, B3}
\seesources
\seealso{T-01-03, T-04-01}
